%
%  chapter1.tex  —  Introdução
%
\chapter{Introdução}
\label{cha:introducao}

\section{Contexto}
\label{sec:contexto}

O Sudoku é um dos puzzles de lógica mais populares do mundo. O seu objectivo é simples: preencher uma grelha $9 \times 9$ com os dígitos de 1 a 9, de tal forma que cada linha, cada coluna e cada uma das nove sub-grelhas $3 \times 3$ (doravante denominadas \emph{caixas}) contenha cada dígito exactamente uma vez. Matematicamente, este problema pertence à classe NP-completo \cite{yato2003complexity}, o que significa que não existe nenhum algoritmo eficiente geral para o resolver. Na prática, no entanto, a estrutura altamente restrita do puzzle permite que algoritmos especializados o resolvam rapidamente.

O Killer Sudoku é uma variante do Sudoku clássico. Para além das restrições de linhas, colunas e caixas, o Killer Sudoku introduz o conceito de \emph{gaiola} (em inglês, \emph{cage}): um grupo de células delimitado por uma linha ponteada, cujos valores devem somar um valor alvo e não se podem repetir dentro da mesma gaiola. Esta restrição adicional torna o espaço de busca simultaneamente mais complexo e mais restrito.

\section{Objectivos}
\label{sec:objectivos}

Este trabalho tem como objectivo o desenvolvimento de quatro solucionadores automáticos para o Sudoku clássico e para o Killer Sudoku, correspondendo a quatro paradigmas algorítmicos distintos:

\begin{enumerate}
    \item \textbf{DFID com MRV} — aprofundamento iterativo com pesquisa em profundidade, implementado em Prolog, com a heurística de valores mínimos restantes (\emph{Minimum Remaining Values}, MRV) para ordenação de variáveis;
    \item \textbf{A*} — pesquisa heurística melhor-primeiro, implementado em Prolog, com uma heurística admissível baseada no número de células vazias;
    \item \textbf{Arrefecimento Simulado (SA)} — algoritmo de optimização metaheurística, implementado em Python, com representação por trocas intra-caixa;
    \item \textbf{Algoritmo Genético (GA)} — algoritmo evolutivo, implementado em Python, com cruzamento preservador de caixas e mutação por troca intra-caixa.
\end{enumerate}

Todos os solucionadores foram estendidos para suportar a variante Killer Sudoku.

\section{Estrutura do Relatório}
\label{sec:estrutura}

O presente relatório está organizado da seguinte forma. O Capítulo~\ref{cha:formulacao} descreve a formulação do problema, a representação adoptada e a implementação detalhada de cada algoritmo. O Capítulo~\ref{cha:resultados} apresenta os resultados experimentais e a análise comparativa dos quatro solucionadores, culminando nas conclusões do trabalho.
